\section*{Chapitre \ref{ch:hf:index-and-basket-arbitrage} --- Arbitrage d'indice et de panier}

\begin{solution}{exo:hf:index-and-basket-arbitrage:1}
$5\,000\,e^{(0.05-0.015)\times0.25}=5\,043.94$~; la bande est comprise entre 5\,043.44 et 5\,044.45.
\end{solution}

\begin{solution}{exo:hf:index-and-basket-arbitrage:2}
Dix contrats micro par E-mini. $1\,000/(5\times150)=1.33$ contrats en dollars par contrat d'Osaka.
\end{solution}

\begin{solution}{exo:hf:index-and-basket-arbitrage:3}
$e^{0.002\times0.25}-1\approx5.0$ points de base.
\end{solution}

\begin{solution}{exo:hf:index-and-basket-arbitrage:4}
$e^{-200/400}=0.61$ et $e^{-600/400}=0.22$.
\end{solution}

\begin{solution}{exo:hf:index-and-basket-arbitrage:5}
Ajouter des noms au-delà du sommet de l'edge réduit l'edge lentement mais le bruit de tracking rapidement (de 0.92 à 0.44 point de base entre cinq et quinze
noms), si bien que le ratio continue de croître jusqu'à ce que la perte de capture et les fourchettes des petits noms dominent.
\end{solution}

\begin{solution}{exo:hf:index-and-basket-arbitrage:6}
Les premières jambes sont les plus susceptibles d'arriver à temps~; y placer les poids les plus lourds maximise le poids capturé.
\end{solution}

\begin{solution}{exo:hf:index-and-basket-arbitrage:7}
Délai moyen de 150 microsecondes~: l'edge culmine à deux noms et l'edge par unité de risque à trois. Délai moyen de 1\,000 microsecondes~: dix et quarante. Plus les
actions suivent, plus de jambes arrivent à temps et plus le meilleur panier est grand.
\end{solution}

\begin{solution}{exo:hf:index-and-basket-arbitrage:8}
Les jambes ne sont pas envoyées d'un seul coup~: dans la simulation du chapitre, la cinquantième jambe arrive une milliseconde après le contrat à terme, alors que les cours affichés des actions ont
déjà bougé. Un backtest qui exécute chaque jambe au prix affiché surestime la capture des grands paniers~; l'edge du panier complet était de 0.07
point de base, et non les 1.35 qu'une exécution instantanée donnerait.
\end{solution}

\begin{solution}{pb:hf:index-and-basket-arbitrage:1}
\begin{enumerate}
\item $F=Se^{(r-q)T}$, et $F(1\pm c)$ avec $c$ les coûts.
\item Le contrat à terme réagit en premier et les cours des actions suivent avec un retard~; les arbitragistes et les teneurs de marché propres aux actions comblent l'écart.
\item La liquidité et la base en valeur absolue se prédisent mutuellement, avec une causalité de Granger dans les deux sens à court terme, sur environ 3\,000 jours.
\item Le trade lent est conservé jusqu'à l'échéance et supporte la prévision de dividende~; le trade rapide est clôturé en quelques secondes.
\item Voir la \cref{def:hf:index-and-basket-arbitrage:partial}~; $(h-w)^\top\Sigma(h-w)$.
\item À chaque étape, on retient le nom qui réduit le plus la variance de tracking, les poids résolvant $\Sigma_{SS}h_S=(\Sigma w)_S$.
\item Voir la \cref{def:hf:index-and-basket-arbitrage:legging}~; $e^{-j\delta/\tau}$.
\item La jambe arrive plus tard, et c'est un nom plus petit avec une fourchette plus large.
\item 1 nom~: 1.44, 0.93, 0.51~; 5~: 2.13, 1.21, 0.92~; 15~: 2.10, 1.39, 0.71~; 50~: 1.72, 1.65, 0.07 (points de base).
\item Edge à 5 noms, edge par unité de risque à 15.
\item L'edge culmine à 10 noms (1.36) et le ratio à 50.
\item Le meilleur panier compte 2 noms à 150 microsecondes et 10 à une milliseconde.
\item À 20 microsecondes par jambe, l'edge net vaut 0.51, 0.92, 0.88, 0.71, 0.32 et 0.07 point de base pour 1, 5, 10, 15, 30 et 50 noms~: cinq noms maximisent
l'edge, et quinze noms l'edge par unité de risque.
\item Voir la \cref{def:hf:index-and-basket-arbitrage:cross}~; le rapport des multiplicateurs multiplié par les devises.
\item Les contrats micro E-mini, un dixième de l'E-mini, depuis le 6 mai 2019~; le Nikkei au SGX et au CME à \textyen500 le point d'indice, et à Osaka à \textyen1\,000.
\item Le \omterm{def:hf:index-and-basket-arbitrage:partial}{panier partiel}~: sa meilleure taille est fixée par la vitesse de la passerelle.
\item Il garde ensemble le contrat à terme et les actions, de sorte que les hedgers, d'un côté comme de l'autre, paient un \omterm{def:hf:fair-value:fair}{prix juste}.
\item Davantage de noms, plus petits~: le meilleur panier compterait encore quelques dizaines de noms, avec davantage de tracking error résiduelle.
\item À partir du délai entre le mouvement du contrat à terme et le premier changement de cours de chaque action, par action et par moment de la journée.
\item La course entre les jambes et les cours des actions, évaluée au regard des fourchettes et du tracking error.
\end{enumerate}
\end{solution}

\begin{solution}{iq:hf:index-and-basket-arbitrage:1}
Les données de la fair value (financement, dividendes jusqu'à l'échéance), les coûts du panier et du contrat à terme, le fait que les actions aient déjà bougé (l'écart peut
reposer sur des données périmées), et la taille disponible.
\iqlookfor{données de la fair value et prix exécutables.}
\end{solution}

\begin{solution}{iq:hf:index-and-basket-arbitrage:2}
On minimise $(h-w)^\top\Sigma(h-w)$ en $h$ avec $h_j=0$ hors de $S$~: le gradient sur $S$ vaut $2(\Sigma(h-w))_S=0$, donc $\Sigma_{SS}h_S=(\Sigma w)_S$~; la
fonction objectif est convexe, donc c'est bien le minimum.
\iqlookfor{conditions du premier ordre sur le sous-ensemble.}
\end{solution}

\begin{solution}{iq:hf:index-and-basket-arbitrage:3}
Sérialisation sur une session et une carte réseau, limitation de débit par session imposée par la bourse, et travail par ordre de la passerelle~; préconstruire les ordres,
les répartir entre sessions et cartes, et les envoyer en parallèle.
\iqlookfor{précalcul et sessions en parallèle.}
\end{solution}

\begin{solution}{iq:hf:index-and-basket-arbitrage:4}
Position longue sur les actions exécutées, adossée à un contrat à terme dimensionné pour tout le panier~: une exposition aux facteurs et le tracking error des noms non exécutés. Redimensionner la
couverture par contrat à terme sur les jambes exécutées, tout de suite, puis décider de compléter ou de déboucler.
\iqlookfor{couvrir ce qui est exécuté, immédiatement.}
\end{solution}

\begin{solution}{iq:hf:index-and-basket-arbitrage:5}
Chaque instrument coté tend à mener pendant les heures de son marché d'origine~: Osaka pendant la journée de Tokyo, Singapour et Chicago lorsque Osaka est fermée~; le mesurer plutôt que
le supposer (chapitre~8).
\iqlookfor{leadership variable dans le temps, mesuré.}
\end{solution}

\begin{solution}{iq:hf:index-and-basket-arbitrage:6}
Capture espérée $\propto\sum_{j=1}^{k}e^{-j\delta/\tau}$ moins $ck$~; il est utile d'ajouter la jambe marginale tant que $e^{-k\delta/\tau}>c$ (en
unités de capture), donc $k^\ast=\lfloor(\tau/\delta)\ln(1/c)\rfloor$.
\iqlookfor{la jambe marginale.}
\end{solution}
